\Qhead{2.5}{Приближение геометрической оптики}{Приближение геометрической оптики. Уравнения эйконала. Дифференциальное уравнение луча. Лучи и поле волны в слоисто-неоднородных средах.}
\Subfile{2.5.1}{sections/02_waves/q05/s01}

\recap{В коротковолновом приближении быстро меняется фаза, медленно — амплитуда. Эйконал задаёт локальный волновой вектор, уравнение переноса — поток; луч подчиняется принципу Ферма.}
\board{Подстановка $A\e^{\ii k_0S}$; два порядка; луч в слоистой среде и поворотная точка.}
\oral{Почему на каустике бесконечна лучевая амплитуда, но не реальное поле?}
\bait{Можно оставить сохранение горизонтальной компоненты волнового вектора: оно сразу даёт закон Снеллиуса для непрерывно слоистой среды.}
\sources{Цысарь~\cite{tsysar1}, лекция~12, PDF-с.~61--63; Сорокин~\cite{sorokinray}, теоретическая подготовка. Уравнения эйконала и луча дополнительно получены прямой подстановкой.}
